\documentclass[10pt,a4paper]{amsart}
\usepackage[margin=19mm]{geometry}
\usepackage{amsmath,amssymb,mathtools}
\usepackage[scr=rsfs,cal=euler]{mathalfa}
\usepackage{xfrac}
\usepackage[unicode,hidelinks,bookmarks=false]{hyperref}
\newcommand{\ie}{i.e.\ }
\newcommand{\cf}{cf.\ }
\newcommand{\resp}{respectively\ }
\newcommand{\Subsection}{\subsection}
\providecommand{\nobreakdash}{\nobreak\hbox{-}}
\setlength{\emergencystretch}{2em}
\multlinegap0pt
\usepackage{bm}
\usepackage{bbm}
\usepackage{dsfont}
\usepackage{xspace}
\usepackage{cancel}
\usepackage{mathtools}
\usepackage{amsthm}
\makeatletter
\@ifundefined{proposition}{\newtheorem{proposition}{Proposition}}{}
\makeatother
\newcommand{\M}{\mathcal{M}}
\newcommand{\R}{\mathbb{R}}
\title[Moduli space of spin connections on three-dimensional homoge]{Moduli space of spin connections on three-dimensional homogeneous spaces}
\author{Matteo Bruno, Gabriele Peluso}
\date{Source: arXiv:2507.06104v2 (CC BY 4.0)}
\begin{document}
\maketitle

\noindent\emph{Source and licence.} Moduli space of spin connections on three-dimensional homogeneous spaces. Version arXiv:2507.06104v2 (CC BY 4.0), distributed under the Creative Commons Attribution 4.0 International licence, as recorded in the arXiv licence metadata for this exact version and shown on the abstract page https://arxiv.org/abs/2507.06104v2. Full rights notice: licenses/arxiv-2507.06104-CC-BY-4.0.txt.\par\smallskip

\noindent\textit{Editorial scope.} Counted unit: the Proposition whose source label is \texttt{None} in the article (source0.tex lines 313--315, source label \texttt{None}), delivered together with the article's own complete original proof of it (source0.tex lines 316--360, one \texttt{proof} environment of 45 lines). No mathematics is written, repaired or re-derived: statement and proof are the article's own text line for line. Required context retained uncounted: none. Every definition, display and label of those lines is retained unchanged. Omitted from this package and not redistributed: 1--312, 361--526. Nothing inside a retained excerpt is omitted or altered: every retained body is delivered exactly as the article's lines are written, so no omission receipt of any kind accompanies this package. Citations: the retained excerpts carry no citation command at all, so no bibliography entry of the article is transcribed and no reference is redistributed. Reference adaptation: none was needed and none was made; every reference inside the delivered unit resolves inside it, so no unresolved reference or citation is produced. Printed slips kept literal: no printed word, symbol, hypothesis or number of the counted statement or of its proof was altered. Layout and class: the article's own class is \texttt{article} with packages including amssymb, lipsum, amsthm, amsfonts, mathrsfs, systeme, mathtools, stmaryrd, tikz-cd, xcolor, graphicx, float, dcolumn, bm; the wrapper is the lane's pinned \texttt{amsart} editorial wrapper at 10pt on A4, which restates the theorem-like environments the retained text needs and adds the article's own macro definitions that the retained text needs (\texttt{\textbackslash M}, \texttt{\textbackslash R}), with extra wrapper packages (bm, bbm, dsfont, xspace, cancel, mathtools). The delivered numbering is the unit's own. Assets: no retained span contains a figure environment, a graphics-inclusion command, a PDF-page inclusion, a table or a TikZ picture; no image file is copied into this package, the package declares no asset dependency and no image file is redistributed with it. Licence: version arXiv:2507.06104v2 of this article is distributed under the Creative Commons Attribution 4.0 International licence, as recorded in the arXiv licence metadata for this exact version and shown on the abstract page https://arxiv.org/abs/2507.06104v2. A full rights notice is delivered at licenses/arxiv-2507.06104-CC-BY-4.0.txt. Characters: every retained excerpt is pure ASCII, so the delivered engine, XeLaTeX with its default Latin Modern text font, reports no missing character.
\begin{proposition}
    There exists a weakly stratified isomorphism between $\mathcal{M}$ and ${\rm Sym}_0(\R^3)\times \R$.
\end{proposition}

\begin{proof}
     To provide the topological manifold structure, we specify a bijective map $\phi:\M\to{\rm Sym}_0(\R^3)\times \R$, which, using the isomorphism ${\rm Sym}_0(\R^3)\times \R\cong \R^6$, gives us a global chart for $\M$. Specifically, we are going to construct a global chart for $\mathcal{C}_{+}$ and $\mathcal{C}_-$, showing that the gluing of them gives us a global chart for the whole $\mathcal{M}$. We define the following maps
    \begin{align*}
        \phi_+:\mathcal{C}_{+}&\to {\rm Sym}_0(\R^3)\times \R_+;\\
        P&\mapsto \left(P-\tfrac{1}{3}tr(P)\mathrm{Id}_3, \lambda_P\right),
    \end{align*}
    where $\lambda_P$ is the eigenvalue of $P$ with the smallest modulus $|\lambda_P|$. This map is bijective with an inverse
    \begin{align*}
        \phi_+^{-1}:{\rm Sym}_0(\R^3)\times \R_+&\to\mathcal{C}_{+};\\
        (A,\lambda)&\mapsto A+(\lambda-\mu_A)\mathrm{Id}_3,
    \end{align*}
    where $A$ is a null-trace symmetric $3\times 3$ matrix and $\mu_A$ is its smallest eigenvalue (surely non-positive). Hence, this map defines a global chart for $\mathcal{C}_{+}$.\\
    With a change of sign, we obtain the global chart for $\mathcal{C}_{-}$
    \[\begin{matrix}
        \phi_-:&\mathcal{C}_{-}&\to &{\rm Sym}_0(\R^3)\times \R_-;\\
        &P&\mapsto &\left(-P+\tfrac{1}{3}tr(P)\mathrm{Id}_3, \lambda_P\right).\\
        & & & \\
        \phi_-^{-1}:&{\rm Sym}_0(\R^3)\times \R_-&\to&\mathcal{C}_{-};\\
        &(A,\lambda)&\mapsto& -A+(\lambda+\mu_A)\mathrm{Id}_3.
    \end{matrix}\]
    Remains to show that the gluing of these two maps is well-defined. Considering $P\in\partial\mathcal{C}_{+}$ and so $-P\in\partial\mathcal{C}_{-}$, hence $\lambda_P=0=\lambda_{-P}$, thus
    \[\phi_-(-P)=\left(-(-P)+\tfrac{1}{3}tr(-P)\mathrm{Id}_3,0\right)=\left(P-\tfrac{1}{3}tr(P)\mathrm{Id}_3,0\right)=\phi_+(P).\]
    We can check that the projection to the quotient $\pi:M(3,\R)\to\mathcal{M}$ is continuous with respect to the global chart $(\mathcal{M},\phi)$. Let us consider the map $M\mapsto\mathrm{sgn}(\det(M))\sqrt{M^tM}$, where $\mathrm{sgn}(\det(M))$ can be $\pm1$ indifferently when $\det(M)$ vanishes, and noticing $\lambda_{-P}=-\lambda_P$ in general, we obtain
\begin{equation}
    \phi\circ \pi(M)=\left(\sqrt{M^tM}-\tfrac{1}{3}tr(\sqrt{M^tM})\mathrm{Id}_3,\,\mathrm{sgn}(\det(M))\lambda_{\sqrt{M^tM}}\right),
\end{equation}
    which is a continuous map.

    \medskip

    We now prove that the strata are mapped appropriately. Evidently, $\phi(M_3)=S_3$ and $\phi(M_0)=S_0$. Moreover, it is immediate to see that $\phi(M_2)\subset S_2$. Conversely, ${\rm rk}(\phi^{-1}(A))\leq 2$, indeed, the eigenvector of $A$ with the smallest eigenvalues $\mu_A$ is in the kernel of $\phi^{-1}(A)=A-\mu_A\mathrm{Id}_3$, then $\phi(M_2)= S_2$.
    
    Let us now consider $P\in M_1$, given a basis of eigenvectors $v_0,v_1,v_2$, with $v_1,v_2\in{ ker}(P)$ and $v_o$ with eigenvalues $\lambda\geq 0$ (without loss in generality), they are eigenvectors of $\phi(P)=P-\frac{1}{3}tr(P)\mathrm{Id}_3$ with eigenvalues $-\frac{\lambda}{3}$ and $\frac{2\lambda}{3}$, respectively. Hence, $\det(\phi(P))=\frac{2}{27}\lambda^3\geq 0$ and ${\rm min}(\phi(P))=-\frac{\lambda^2}{3}$, from which $27(\det(A))^2=-4({\rm min}(A))^3$. Then, $\phi(M_1)\subset S_1$\\
    Let $A\in S_1$, its eigenvalues $\mu_1,\mu_2,\mu_3$ satisfy
    \begin{equation*}
        \mu_1+\mu_2+\mu_3=0,\,\text{ and }\ \tfrac{1}{4}(\mu_1\mu_2\mu_3)^2=-\tfrac{1}{27}(\mu_1\mu_2+\mu_2\mu_3+\mu_3\mu_1)^3
    \end{equation*}
    Plugging the first into the second, we obtain an equation
    \[4\mu_1^6+12\mu_1^5\mu_2-3\mu_1^4\mu_2^2-26\mu_1^3\mu_2^3-3\mu_1^2\mu_2^4+12\mu_1\mu_2^5+4\mu_2^6=0.\]
    A possible solution is $\mu_2=0$, in this case $\det(A)=0$, and so ${\rm min}(A)=0$ and $tr(A)=0$, imposing that $A=0$. Considering $\mu_2\neq 0$, we can introduce the variable $t=\mu_1/\mu_2$, and the previous equation reads
    \[4t^6+12t^5-3t^4-26t^3-3t^2+12t+4=0.\]
    This polynomial equation has three roots, each with multiplicity $2$: $t=1$, $t=-2$, $t=-\frac{1}{2}$, providing three similar cases: $\mu_1=\mu_2=-\frac{1}{2}\mu_3$, $\mu_2=\mu_3=-\frac{1}{2}\mu_1$, and $\mu_3=\mu_1=-\frac{1}{2}\mu_2$. The condition $\det(A)>0$ forces the two equal eigenvalues to be negative. Thus, the spectrum is 
    \[\sigma(A)=\{-\mu,-\mu,2\mu\}\ \text{ with }\ \mu\geq 0.\]
    Now, $\phi^{-1}(A)$ is given by $A+\mu\mathrm{Id}_3$ which has rank almost $1$. Thus, $\phi(M_1)=S_1$.
\end{proof}

\end{document}
